\documentclass[11pt,a4paper]{ltjsarticle}
% Compile with LuaLaTeX
% If you use pLaTeX/dvipdfmx, change the documentclass if needed.

% =========================
% Packages
% =========================
\usepackage{luatexja}
\usepackage{graphicx}
\usepackage{hyperref}
\usepackage{amsmath,amssymb,mathtools}
\usepackage{xcolor}
\usepackage{float}
\usepackage{subcaption}
\usepackage{booktabs}
\usepackage{enumitem}
\usepackage{geometry}

\geometry{margin=25mm}
\setcounter{tocdepth}{3}

% =========================
% User information
% =========================
\title{Weekly Report: Project Title}
\author{Your Name}
\date{Week: YYYY/MM/DD -- YYYY/MM/DD}

% =========================
% Custom commands
% =========================
\newcommand{\sectionblock}[1]{%
  \vspace{4mm}
  {\large\textbf{#1}}
  \vspace{2mm}

}

\newcommand{\todo}[1]{\textcolor{red}{#1}}

\begin{document}

\maketitle

% Optional table of contents
% \tableofcontents
% \newpage



% =========================================================
% 1. Question / Purpose
% =========================================================
\sectionblock{Question}

% What is the scientific or technical question this week?
% Keep it concrete.

\begin{quote}
What do you want to understand, demonstrate, or improve this week?
\end{quote}

% =========================================================
% 2. Hypothesis
% =========================================================
\sectionblock{Hypothesis}

% Describe your expectation before the experiment/simulation.
% Example:
% If the applied pressure generates a transverse strain component opposite to the intrinsic strain,
% the ODMR splitting should become smaller.


% \begin{figure}[H]
%   \centering
%   \includegraphics[width=0.75\linewidth]{image/example_result.png}
%   \caption{
%   Example result.
%   Replace this figure with the main ODMR spectrum, pressure estimation, force feedback result, or simulation result.
%   }
%   \label{fig:main_result}
% \end{figure}



% =========================================================
% 3. Evidence / Interpretation
% =========================================================
\sectionblock{Evidence / Interpretation}

% Explain why the result supports or does not support your hypothesis.
% Use equations if necessary.




% Add your interpretation here:
\begin{itemize}[leftmargin=8mm]
  \item Evidence 1:
  \item Evidence 2:
  \item Interpretation:
\end{itemize}

% =========================================================
% 4. Gap / Failure
% =========================================================
\sectionblock{Gap / Failure}

% Be honest about what did not work.
% This section is very useful for discussion with collaborators/supervisors.

\begin{itemize}[leftmargin=8mm]
  \item Possible reason:
  \item What is still unclear:
\end{itemize}

% =========================================================
% 5. Next Step
% =========================================================
\sectionblock{Next Step}



% =========================================================
% 6. References
% =========================================================
\sectionblock{References}

% Example:
% [1] Author Name, Journal Name, Volume, Page (Year).
% [2] Author Name, arXiv:XXXX.XXXXX.

\begin{enumerate}[leftmargin=8mm]
  \item Reference 1
  \item Reference 2
\end{enumerate}

\end{document}
